\section*{Author contributions}
G.J.W.\ conceived the survey, defined the target sample and the selection
criteria, and wrote the manuscript. R.D.\ built and operated the archival
retrieval, calibration and search pipeline, implemented the spatial control
statistic and the injection--recovery campaigns, and produced the per-window
catalogue and its regeneration scripts. Both authors designed the vetting
procedure, chose the masked species, carried out the statistical
validation reported in \S\ref{sec:statistic} and
Appendix~\ref{app:falsealarm}, and agreed the dispositions of
Table~\ref{tab:ledger}.
%% AUTHORS: confirm this division before submission, and see
%% AUTHOR_ACTIONS.md item 6 for the affiliation statement this section
%% should be read with.

\section*{Funding}
This work received no dedicated funding. It used public archival data and
no observing time was awarded for it.
%% AUTHORS: replace if any grant supported this work.

\section*{Competing interests}
The authors declare no competing interests. VBRL Holdings is a privately
held company through which R.D.\ conducts independent research, named
here solely as his affiliation; it had no role in the design of the
survey, the selection of targets, the analysis, the decision to publish
or the content of this paper, and provided no funding for the work.
%% AUTHORS: confirm this wording, which is written to be accurate only if
%% the company's involvement was indeed nil. See AUTHOR_ACTIONS.md item 6.

\section*{Acknowledgements}
This paper uses ALMA data from the public ALMA Science Archive. The \NEB{}
primary-census execution blocks searched span \PrEbDateFirst{} to
\PrEbDateLast{} and belong to \NProjCodes{} ALMA proposals, listed in
Table~\ref{tab:projcodes}, whose principal investigators and teams we
thank.\ProjCodeUnresClause{} ALMA is a partnership of ESO
(representing its member states), NSF (USA) and NINS (Japan), together with NRC
(Canada), NSTC and ASIAA (Taiwan), and KASI (Republic of Korea), in cooperation
with the Republic of Chile. The Joint ALMA Observatory is operated by ESO,
AUI/NRAO and NAOJ. This work has made use of data from the European Space
Agency (ESA) mission {\it Gaia} \citep{GaiaDR3}, processed by the {\it Gaia}
Data Processing and Analysis Consortium (DPAC). This research has made use
of the SIMBAD database \citep{SIMBAD2000} (CDS, Strasbourg, France), of the NASA Exoplanet
Archive \citep{NEA2013}, and of the Splatalogue spectral-line database
\citep{Splatalogue2007}.

\section*{Data Availability}
\label{sec:data}
Every observation this paper rests on is public. The \NEB{} primary-census
execution blocks and the blocks of every other set named in
Table~\ref{tab:blockledger} are held in the ALMA Science Archive and are
retrieved there by the execution-block identifier, which
Tables~\ref{tab:ledger} and~\ref{tab:ledgercont} give for every one of the
\EvNCrossRaw{} threshold crossings. The ancillary catalogues are public in the
same sense: Gaia~DR3 for the astrometry and the colours, the NASA Exoplanet
Archive for the planet counts, and the CDMS and JPL line catalogues through
Splatalogue for the mask, queried once on \McQueryDate.

The derived results are in the paper. Table~\ref{tab:blockledger} partitions
the blocks, windows, stars, systems and hours of every set the search defines;
Table~\ref{tab:counts} carries every headline count, the five fates of the
unprocessed blocks and the crossing partition; Tables~\ref{tab:ledger}
and~\ref{tab:ledgercont} list all \EvNCrossRaw{} crossings with their
frequencies in both frames, the stellar and control statistics, the nearest
masked transition and the recurrence test; Table~\ref{tab:cand} gives the two
leading events in full; and Table~\ref{tab:holdout} does the same for the
reserved hold-out. The method is set out in enough detail to be
reimplemented, and no separate data deposit accompanies the paper.

Three choices that could create or remove an unattributed crossing carry an
identifier anyone can check. The
calibration hold-out rule of Appendix~\ref{sec:heldout} and the block list
its first guard depends on were committed to the project repository as
\texttt{c75069040eab52ec1336376a4dae62114bb281a3} on 2026 September~14 at
07:10~UTC, \HoLeadBlockDays{} days before the first reserved block was
searched. The molecular-line list is a single dated query of a public line
catalogue (\McQueryDate), so it is specified rather than
described, and it post-dates the fixing of the disposition criteria on
\MkCritDate{} (\texttt{f5f52e7804f1baad99515ea2d46ddd7e4e4cd546}), as does
every species in it, so no species in the list was recognised late or early
with respect to them.
%% r15-deposit, 2026-10-08.  GLENN'S DECISION, VERBATIM: "For the Zenodo
%% deposit, remove all traces of this offer from the paper, I will not be
%% doing it."  This statement therefore promises nothing.  It names only what
%% a reader can actually obtain: the ALMA archive, the public ancillary
%% catalogues, the tables in this paper, and the two commits -- both of which
%% resolve publicly (checked 2026-10-08; an all-zero control SHA 404s).
%% `doigate.py` is INVERTED: it now FAILS if an unfulfilled promise of a
%% deposit reappears anywhere in the manuscript.  Do not reintroduce
%% "available on request" or "will be deposited" without the author's word.
%% Referee 2's minor 14 ("insert the DOI") is DECLINED on the author's
%% instruction, and the response letter must say so plainly.

\section*{Software}
The pipeline depends on CASA~6.7.6 \citep[casatools/casatasks
6.7.6-14;][]{CASA2022}, Python~3 with astropy \citep{Astropy2022}, numpy
\citep{Harris2020} and scipy \citep{Virtanen2020}, and the ALMA Science
Archive and SIMBAD TAP interfaces through astroquery \citep{Ginsburg2019};
all are free and open-source, at the versions cited.
